\section{(Un)Soundness of Method Local Classes}
With our language primitives in place, we now consider type safety of two variants of our language.
The two variants differ only in the typing of local classes:
\begin{enumerate}
\item $\zeta(\mathsf{Java})$ types local classes with rule \textsc{new-local-java} of Section~\ref{sec:funnelling}, as Java does;
\item $\zeta(\mathsf{funnelled})$ types local classes with rule \textsc{new-local} of Section~\ref{sec:typing}.
\end{enumerate}
Both variants use the type test of Java:
rule \textsc{if-instanceof} accepts only checked casts (Section~\ref{sec:cast}), and at run time rules \textsc{ifThen} and \textsc{ifElse} compare class names only.
As shown in Section~\ref{sec:core}, $\zeta(\mathsf{Java})$ is unsound.
We now show that $\zeta(\mathsf{funnelled})$ is sound.
In the statements below, the program $\decls\;e$ satisfies well-formedness requirements A1-A9, and $\mathsf{Dec}~\mathrm{Ok}$ holds for every $\mathsf{Dec}\in\decls$.

\begin{lemma}[Well-Formedness]
If $\step{e}{e'}$, then the program $\decls\;e'$ satisfies well-formedness requirements A1-A9.
\end{lemma}

\begin{theorem}[Progress]
If $\typing{\emptyset}{\emptyset}{e}{T}$, then either $e$ is a value $v$, or there is an $e'$ such that $\step{e}{e'}$.
\end{theorem}

\begin{theorem}[Preservation]
If $\typing{\emptyset}{\emptyset}{e}{T}$ and $\step{e}{e'}$, then $\typing{\emptyset}{\emptyset}{e'}{T}$.
\end{theorem}

\begin{corollary}[Type Soundness]
If $\vdash\decls\;e_0~\mathrm{Ok}$ with $\typing{\emptyset}{\emptyset}{e_0}{T}$, and $e_0\rightarrow^* e$ with no $e'$ such that $\step{e}{e'}$, then $e$ is a value $v$ and $\typing{\emptyset}{\emptyset}{v}{T}$.
\end{corollary}
% Reified type tests are left out of the formal model on purpose; see Section 2 and the comment after rule ifElse in Figure fig:opsem.
